\chapter{Dirichlet's Theorem}

\begin{center}\itshape
Every unbounded arithmetic progression whose first term and common difference are integers without common factor
contains infinitely many primes.\\
\upshape --- Title of Dirichlet's memoir (1837)
\end{center}

In 1853 Chebyshev wrote to his colleague Fuss that something odd seemed to be going on: there appeared to be more primes of the form $4n+3$ than of
the form $4n+1$. Is the race rigged? Does one team eventually run away with all the primes? This chapter proves that, in the long run, every
team gets its fair share; Chapter~7 explains why the race nevertheless looks rigged for as far as anyone can see.

Euclid's argument adapts to a few progressions: $4n+3$ (Tome~1, Exercise~2.1), $6n+5$, and with more ingenuity $kn+1$ for every $k$
(Tome~1, Exercise~5.7). For $5n+2$ no such trick is known, and there are reasons to think none exists for most progressions. Dirichlet's
analytic method treats all of them at once. With the $L$-functions of Chapter~3 and their non-vanishing at $1$, what remains is
bookkeeping, and the result is quantitative: in the weighted sense of Mertens's first theorem, the primes are \emph{equally
distributed} among the $\varphi(q)$ classes prime to $q$.

\begin{think}
Up to $x=60\,000$ there are at least as many primes $\equiv3\pmod4$ as primes $\equiv1\pmod4$ at every $x$, except on the tiny interval $26861\le x<26863$ (Figure~\ref{fig:race}). Does that
contradict equal distribution? What would ``equally distributed'' have to mean for such a race to be compatible with it?
\end{think}

\section{The theorem}

Throughout, $q\ge1$ is fixed, $a$ is an integer prime to $q$, and the constants implied by $O(\cdot)$ may depend on $q$.

\begin{theorem}[Dirichlet, 1837]\label{thm:dirichlet}
For $x\ge2$,
\[
\sum_{\substack{p\le x\\p\equiv a\,(q)}}\frac{\log p}p=\frac1{\varphi(q)}\log x+O(1).
\]
In particular, there are infinitely many primes $p\equiv a\pmod q$.
\end{theorem}

Let $S(x,\chi)=\sum_{n\le x}\frac{\chi(n)\Lambda(n)}n$. The proof filters the class of $a$ with characters (Tome~2, Proposition~11.15):
\[
\sum_{\substack{n\le x\\n\equiv a\,(q)}}\frac{\Lambda(n)}n=\frac1{\varphi(q)}\sum_\chi\overline{\chi(a)}\,S(x,\chi).
\]
The principal character gives the main term, and each other character contributes $O(1)$.

\begin{lemma}\label{lem:principal}
$S(x,\chi_0)=\log x+O(1)$.
\end{lemma}
\begin{proof}
$S(x,\chi_0)$ differs from $\sum_{n\le x}\frac{\Lambda(n)}n=\log x+O(1)$ (Theorem~\ref{thm:mertens1}) by the terms with $n=p^k$, $p\mid q$, which add up to at most
$\sum_{p\mid q}\log p\sum_{k\ge1}p^{-k}=\sum_{p\mid q}\frac{\log p}{p-1}$, a constant.
\end{proof}

\begin{lemma}\label{lem:chisums}
Let $\chi\ne\chi_0$. For $y\ge1$:
\[
\begin{gathered}\sum_{m\le y}\frac{\chi(m)}m=L(1,\chi)+O\Bigl(\frac1y\Bigr),\qq\sum_{m\le y}\frac{\chi(m)\log m}m=O(1).\end{gathered}
\]
\end{lemma}
\begin{proof}
Let $A(t)=\sum_{n\le t}\chi(n)$, so $|A|\le\varphi(q)$ (Proposition~\ref{prop:orth}). For $f$ of class $C^1$, decreasing and tending to $0$ on $[y,\infty[$, Lemma~\ref{lem:abelint}
gives $\sum_{y<n\le x}\chi(n)f(n)=A(x)f(x)-A(y)f(y)-\int_y^xA(t)f'(t)\dd t$, so $\bigl|\sum_{n>y}\chi(n)f(n)\bigr|\le2\varphi(q)f(y)$. With $f(t)=\frac1t$, the tail of $L(1,\chi)$ after $y$ is
at most $\frac{2\varphi(q)}y$. With $f(t)=\frac{\log t}t$, decreasing on $[3,\infty[$, the series $\sum\frac{\chi(m)\log m}m$ converges (its tail after $y\ge3$ is at most $2\varphi(q)\frac{\log y}y$),
so its partial sums are bounded.
\end{proof}

\begin{proposition}\label{prop:nonprincipal}
Let $\chi\ne\chi_0$. Then $S(x,\chi)=O(1)$.
\end{proposition}
\begin{proof}
Since $\log=\Lambda*\mathbf1$ (Chapter~1) and $\chi$ is completely multiplicative, $\chi(n)\frac{\log n}n=\sum_{dm=n}\frac{\chi(d)\Lambda(d)}d\cdot\frac{\chi(m)}m$. Sum over $n\le x$ and group
by $d$ (Lemma~\ref{lem:convsum}):
\[
\begin{aligned}\sum_{n\le x}\frac{\chi(n)\log n}n&=\sum_{d\le x}\frac{\chi(d)\Lambda(d)}d\sum_{m\le x/d}\frac{\chi(m)}m\eqbr=\sum_{d\le x}\frac{\chi(d)\Lambda(d)}d\bigl(L(1,\chi)+O(\tfrac dx)\bigr)\end{aligned}
\]
by Lemma~\ref{lem:chisums}. The error terms add up to $O\bigl(\frac{\psi(x)}x\bigr)$, which is $O(1)$ by Proposition~\ref{prop:chebyshev}(a), and the left
side is $O(1)$ by Lemma~\ref{lem:chisums}. So $L(1,\chi)S(x,\chi)=O(1)$, and since $L(1,\chi)\ne0$ (Theorem~\ref{thm:lnonzero}), $S(x,\chi)=O(1)$.
\end{proof}

\begin{proof}[Proof of Theorem~\ref{thm:dirichlet}]
By the filter, Lemma~\ref{lem:principal} and Proposition~\ref{prop:nonprincipal}, using $\chi_0(a)=1$ and $|\chi(a)|=1$,
\[
\sum_{\substack{n\le x\\n\equiv a\,(q)}}\frac{\Lambda(n)}n=\frac1{\varphi(q)}\log x+O(1).
\]
The prime powers $p^k$ with $k\ge2$ contribute at most $\sum_p\frac{\log p}{p(p-1)}<\infty$ (as in the proof of Theorem~\ref{thm:mertens1}), which gives the formula.
Its right side tends to infinity, so the sum on the left has infinitely many terms.
\end{proof}

\begin{beauty}
Look at where $L(1,\chi)\ne0$ entered: the identity $\sum\chi(n)\frac{\log n}n=L(1,\chi)S(x,\chi)+O(1)$ says that the weighted primes twisted by $\chi$ form a
bounded sum \emph{unless} $L(1,\chi)=0$. If $L(1,\chi)$ vanished, the same identity with one more logarithm would force $S(x,\chi)\sim-\log x$ (Exercise~\ref{ex:ifzero}):
the primes would conspire to avoid the classes where $\chi$ is $1$ and crowd into the others. For a real $\chi$ modulo $4$, half the classes would get
no primes at all, in the logarithmic sense. Non-vanishing at $s=1$ is equidistribution, and both statements are about the same object.
\end{beauty}

\section{Consequences}

With the theorem in hand, Mertens's three results of Chapter~1 can be split class by class at no extra cost.

\begin{corollary}\label{cor:recipprog}
There is a constant $M(q,a)$ such that, for $x\ge2$,
\[
\sum_{\substack{p\le x\\p\equiv a\,(q)}}\frac1p=\frac1{\varphi(q)}\log\log x+M(q,a)+O\Bigl(\frac1{\log x}\Bigr).
\]
\end{corollary}
\begin{proof}
Write $\sum_{p\le t,\,p\equiv a}\frac{\log p}p=\frac1{\varphi(q)}\log t+R_a(t)$ with $R_a$ bounded, and repeat the proof of Theorem~\ref{thm:mertens2} with $f(t)=\frac1{\log t}$: the main
term gives $\frac1{\varphi(q)}(\log\log x-\log\log2)$ and the integral of $\frac{R_a(t)}{t\log^2t}$ converges, with a tail $O(\frac1{\log x})$.
\end{proof}

Summing over the classes, $\sum_aM(q,a)=M-\sum_{p\mid q}\frac1p$, since the classes not prime to $q$ contain only the primes dividing $q$ (Exercise~\ref{ex:constmod4}).
Figure~\ref{fig:mertensprog} shows the convergence for $q=3$.

\begin{figure}[ht]
\centering
\fig{ch04-mertensprog}
\caption{$\sum_{p\le x,\,p\equiv a\,(3)}\frac{\log p}p-\frac12\log x$ for $a=1$ (blue) and $a=2$ (orange), $10\le x\le10^6$, in logarithmic scale. Both are bounded, as
Theorem~\ref{thm:dirichlet} says, and both seem to converge, to about $-1.187$ and $-0.512$. The gap, about $0.675$, is not a bias of large primes: about $0.368$ of it is
$\frac{L'}L(1,\chi_{-3})$, which reflects the whole sequence of primes (the prime $2\equiv2$ among them), and $0.308$ comes from the prime powers $p^k\equiv1$, $k\ge2$, which a sum
over primes leaves out. Chapter~5 shows that the limits exist.}
\label{fig:mertensprog}
\end{figure}

\section{Equal distribution, and in what sense}

Theorem~\ref{thm:dirichlet} gives each class prime to $q$ the same share $\frac1{\varphi(q)}$ of the \emph{logarithmically weighted} primes. Counting primes
without weights needs more. Let $\pi(x;q,a)$ be the number of primes $p\le x$ with $p\equiv a\pmod q$. The \emph{prime number theorem for arithmetic progressions}
says that
\[
\pi(x;q,a)\sim\frac1{\varphi(q)}\frac x{\log x},
\]
which follows from the method of Chapter~5 applied to $L(s,\chi)$ on the line $\operatorname{Re}s=1$; Chapter~5 indicates the changes. Even that says nothing
about which class leads at a given $x$, and Figure~\ref{fig:race} shows that the leader can be remarkably stable.

\begin{figure}[ht]
\centering
\fig{ch04-race}
\caption{The race $\pi(x;4,3)-\pi(x;4,1)$ for $x\le60\,000$. The class $3$ is ahead or level everywhere (level at only ten primes, from $5$ to $26921$), except at $x=26861$,
where class $1$ takes the lead by one, until the next prime $26863\equiv3$. In the long run the class $3$ leads for about
$99.59\%$ of the time, in the logarithmic sense: this is \emph{Chebyshev's bias}, explained in Chapter~7.}
\label{fig:race}
\end{figure}

\begin{remark}
The bias has a simple heuristic source. The squares are all $\equiv1\pmod4$ (for odd numbers), and $\psi(x;4,1)$ counts the squares of primes,
while $\pi$ does not. Passing from the smooth weights of $\psi$ to the raw count $\pi$ subtracts roughly $\frac{\sqrt x}{\log x}$ from class $1$; the error terms
of the prime number theorem in progressions are of the same size under the Riemann hypothesis for $L(s,\chi_{-4})$, so the subtraction is
visible. Rubinstein and Sarnak (1994) made this precise, conditionally; the figure $99.59\%$ is theirs. Chapter~7 develops this.
\end{remark}

\section{Exercises}

\begin{exercise}[$\star$]\label{ex:euclid}
Show that there are infinitely many primes $p\equiv5\pmod6$ and infinitely many primes $p\equiv1\pmod4$, by Euclid-style arguments.
(For the second, consider the prime factors of $(2p_1\cdots p_k)^2+1$.)
\end{exercise}

\begin{exercise}[$\star$]\label{ex:constmod4}
Show that $M(4,1)+M(4,3)=M-\frac12$, with $M$ the Meissel--Mertens constant, and more generally $\sum_{a}M(q,a)=M-\sum_{p\mid q}\frac1p$, the sum
running over the classes prime to $q$.
\end{exercise}

\begin{exercise}[$\star\star$]\label{ex:density}
\emph{Dirichlet density.} Show that $\sum_{p\equiv a\,(q)}p^{-\sigma}=\frac1{\varphi(q)}\log\frac1{\sigma-1}+O(1)$ as $\sigma\to1^+$, directly from Theorem~\ref{thm:lbasic} and
Theorem~\ref{thm:lnonzero}. This was Dirichlet's original formulation.
\end{exercise}

\begin{exercise}[$\star\star$]\label{ex:progproduct}
Show that there is a constant $C(q,a)>0$ with $\prod_{p\le x,\,p\equiv a\,(q)}\bigl(1-\frac1p\bigr)\sim C(q,a)(\log x)^{-1/\varphi(q)}$.
\end{exercise}

\begin{exercise}[$\star\star$]\label{ex:two}
\emph{Signature exercise.} Show that the primes $p$ for which $2$ is a quadratic residue modulo $p$ get half of the weight:
$\sum_{p\le x,\,(2/p)=1}\frac{\log p}p=\frac12\log x+O(1)$. Do the same for $-1$ and $-3$ in place of $2$.
\end{exercise}

\begin{exercise}[$\star\star$]\label{ex:ifzero}
Suppose, contrary to Theorem~\ref{thm:lnonzero}, that $L(1,\chi)=0$ for some $\chi\ne\chi_0$, and let $M_\chi(x)=\sum_{n\le x}\frac{\mu(n)\chi(n)}n$.
\begin{enumerate}[label=(\alph*)]
\item Show that $L'(1,\chi)M_\chi(x)=\log x+O(1)$.
\item Show that $S(x,\chi)=-L'(1,\chi)M_\chi(x)+O(1)$, and conclude that $S(x,\chi)=-\log x+O(1)$.
\end{enumerate}
\end{exercise}

\section{Hints}

\begin{hint}{ex:euclid}
For $5\pmod6$: consider $6p_1\cdots p_k-1$. For $1\pmod4$: if $p\mid m^2+1$ with $p$ odd, then $-1$ is a square modulo $p$.
\end{hint}
\begin{hint}{ex:constmod4}
Add the asymptotic formulas of Corollary~\ref{cor:recipprog} over $a$ and compare with Theorem~\ref{thm:mertens2}.
\end{hint}
\begin{hint}{ex:density}
Filter with characters. Use $\log L(\sigma,\chi)=\sum_p\chi(p)p^{-\sigma}+O(1)$, $\log L(\sigma,\chi_0)=\log\frac1{\sigma-1}+O(1)$, and that $L(\sigma,\chi)$ tends to a non-zero limit for $\chi\ne\chi_0$.
\end{hint}
\begin{hint}{ex:progproduct}
Take logarithms; $\log(1-\frac1p)=-\frac1p+O(\frac1{p^2})$.
\end{hint}
\begin{hint}{ex:two}
$\bigl(\frac2p\bigr)=\chi_8(p)$, the real character modulo $8$ with $\chi_8(n)=1$ for $n\equiv\pm1\pmod8$ (Tome~1, Chapter~6). The set of primes is the union of two
classes modulo $8$.
\end{hint}
\begin{hint}{ex:ifzero}
(a) Evaluate in two ways
\[
\sum_{n\le x}\frac{\mu(n)\chi(n)}n\sum_{m\le x/n}\frac{\chi(m)}m\log\frac x{nm}.
\]
(b) $\Lambda=\mu*\log$. In both, use Lemma~\ref{lem:chisums}, Exercise~\ref{ex:lderivative} and
\[
\sum_{m\le y}\frac{\chi(m)\log m}m=-L'(1,\chi)+O\Bigl(\frac{1+\log y}y\Bigr).
\]
\end{hint}

\section{Solutions}

\begin{solution}{ex:euclid}
\idea A number of the shape $6P-1$ must have a prime factor $\equiv5\pmod6$, since a product of primes $\equiv1\pmod6$ is $\equiv1$. For $1\pmod4$, the
prime factors of $m^2+1$ are forced into the class $1$ by quadratic residues.
\writeup Suppose $p_1,\dots,p_k$ are all the primes $\equiv5\pmod6$, and let $N=6p_1\cdots p_k-1$. Its prime factors are not $2$ or $3$, so they are $\equiv\pm1\pmod6$;
they are not all $\equiv1$, because $N\equiv5$; and none is a $p_i$, since $p_i\mid N+1$. Contradiction. For $1\pmod4$, let $N=(2p_1\cdots p_k)^2+1$, odd and
$>1$. A prime $p\mid N$ is odd, not a $p_i$, and $-1\equiv(2p_1\cdots p_k)^2$ is a square modulo $p$, so $p\equiv1\pmod4$ by Euler's criterion $\bigl(\frac{-1}p\bigr)=(-1)^{(p-1)/2}$.
These tricks need a polynomial whose prime divisors lie in the target class; such polynomials exist exactly for the classes $a$ with $a^2\equiv1\pmod q$
(a theorem of Murty, 1988), so $5n+2$ is out of reach of Euclid.
\end{solution}

\begin{solution}{ex:constmod4}
\idea The classes partition the primes, apart from the finitely many dividing $q$.
\writeup Summing Corollary~\ref{cor:recipprog} over the $\varphi(q)$ classes prime to $q$ gives $\sum_{p\le x,\,p\nmid q}\frac1p=\log\log x+\sum_aM(q,a)+O(\frac1{\log x})$. The left side is
$\sum_{p\le x}\frac1p-\sum_{p\mid q}\frac1p$ for $x\ge q$, which by Theorem~\ref{thm:mertens2} is $\log\log x+M-\sum_{p\mid q}\frac1p+O(\frac1{\log x})$. Compare the constants. For $q=4$
the only prime dividing $q$ is $2$.
\end{solution}

\begin{solution}{ex:density}
\idea This is Theorem~\ref{thm:dirichlet} in Dirichlet's own language: instead of cutting at $x$, weight by $p^{-\sigma}$ and let $\sigma\to1$. The logarithm of $L(\sigma,\chi)$
is essentially $\sum_p\chi(p)p^{-\sigma}$.
\writeup By Theorem~\ref{thm:lbasic}(a), $\log L(\sigma,\chi)=\sum_p\chi(p)p^{-\sigma}+E(\sigma,\chi)$, where $\log$ denotes the series of that theorem and
$|E(\sigma,\chi)|\le\sum_p\sum_{k\ge2}p^{-k}\le\sum_p\frac1{p(p-1)}<1$. Filtering with characters,
\[
\sum_{p\equiv a}p^{-\sigma}=\frac1{\varphi(q)}\sum_\chi\overline{\chi(a)}\bigl(\log L(\sigma,\chi)-E(\sigma,\chi)\bigr).
\]
For $\chi\ne\chi_0$, $L(\sigma,\chi)\to L(1,\chi)\ne0$, and $\exp(\log L(\sigma,\chi))=L(\sigma,\chi)$, so $\log L(\sigma,\chi)$ is a continuous logarithm of a function with non-zero limit:
it stays bounded as $\sigma\to1^+$. For $\chi_0$, $L(\sigma,\chi_0)=\zeta(\sigma)\prod_{p\mid q}(1-p^{-\sigma})$ is real positive and $\sim\frac{\varphi(q)}q\frac1{\sigma-1}$, and $\log L(\sigma,\chi_0)$ is real
(real coefficients), so it equals $\log\frac1{\sigma-1}+O(1)$. Hence the claim. Dirichlet stated his theorem this way.
\end{solution}

\begin{solution}{ex:progproduct}
\idea As in Mertens's third theorem, the logarithm of the product is minus the sum of $\frac1p$ plus a convergent series; but now there is no
analogue of $\gamma$ to identify the constant.
\writeup With $h$ as in the proof of Theorem~\ref{thm:mertens3},
\[
\log\prod_{\substack{p\le x\\p\equiv a}}\Bigl(1-\frac1p\Bigr)=-\sum_{\substack{p\le x\\p\equiv a}}\frac1p-\sum_{\substack{p\le x\\p\equiv a}}h(p),
\]
where $0\le h(p)=-\log(1-\frac1p)-\frac1p\le\frac1{p^2}$. The second sum converges to some $H(q,a)$ with tail $O(\frac1x)$. With Corollary~\ref{cor:recipprog}, the logarithm is $-\frac1{\varphi(q)}\log\log x-M(q,a)-H(q,a)+O(\frac1{\log x})$.
Exponentiate: $C(q,a)=e^{-M(q,a)-H(q,a)}$. Unlike $e^{-\gamma}$, these constants have no closed form in general; they are computed in the notebook
for $q=4$.
\end{solution}

\begin{solution}{ex:two}
\idea Being a quadratic residue is detected by a real character, and a real non-principal character contributes $O(1)$ by Proposition~\ref{prop:nonprincipal}. So
the primes split evenly between residues and non-residues.
\writeup For odd $p$, $\bigl(\frac2p\bigr)=1$ iff $p\equiv\pm1\pmod8$, that is $\frac{1+\chi_8(p)}2=1$. So
\[
\begin{aligned}\sum_{\substack{p\le x\\(2/p)=1}}\frac{\log p}p&=\frac12\sum_{2<p\le x}\frac{\log p}p+\frac12\sum_{p\le x}\frac{\chi_8(p)\log p}p\eqbr=\frac12\log x+O(1),\end{aligned}
\]
by Theorem~\ref{thm:mertens1} and Proposition~\ref{prop:nonprincipal} (the prime powers again contribute $O(1)$). For $-1$: $\bigl(\frac{-1}p\bigr)=\chi_{-4}(p)$ for odd $p$. For $-3$:
by quadratic reciprocity, $\bigl(\frac{-3}p\bigr)=\bigl(\frac{-1}p\bigr)\bigl(\frac3p\bigr)=\bigl(\frac p3\bigr)=\chi_{-3}(p)$ for $p>3$. The same computation applies. In general, for a non-square $b$ the map
$p\mapsto\bigl(\frac bp\bigr)$ is a non-principal real character modulo $4|b|$ evaluated at $p$, and half of the primes have $b$ as a residue; the bookkeeping with
reciprocity is in Ireland and Rosen, Chapter~5.
\end{solution}

\begin{solution}{ex:ifzero}
\idea If $L(1,\chi)=0$, the derivative $L'(1,\chi)$ takes over the role that $L(1,\chi)$ played in Proposition~\ref{prop:nonprincipal}, and the convolution identities produce a
main term $\log x$ where there used to be nothing.
\writeup Let $y\ge1$. By Lemma~\ref{lem:chisums} and its proof, $\sum_{m\le y}\frac{\chi(m)}m=O(\frac1y)$ (as $L(1,\chi)=0$) and $\sum_{m\le y}\frac{\chi(m)\log m}m=-L'(1,\chi)+O(\frac{1+\log y}y)$, using
$L'(1,\chi)=-\sum_m\frac{\chi(m)\log m}m$ (Exercise~\ref{ex:lderivative}). Hence
\[
\begin{aligned}\sum_{m\le y}\frac{\chi(m)}m\log\frac ym&=\log y\sum_{m\le y}\frac{\chi(m)}m\eqbr-\sum_{m\le y}\frac{\chi(m)\log m}m\eqbr=L'(1,\chi)+O\Bigl(\frac{1+\log y}y\Bigr).\end{aligned}
\]
(a) Let $T=\sum_{n\le x}\frac{\mu(n)\chi(n)}n\sum_{m\le x/n}\frac{\chi(m)}m\log\frac x{nm}$. Grouping by $k=nm$, the coefficient of $\frac1k\log\frac xk$ is $\sum_{nm=k}\mu(n)\chi(n)\chi(m)=\chi(k)\sum_{n\mid k}\mu(n)$,
which is $1$ for $k=1$ and $0$ otherwise; so $T=\log x$. On the other hand, with $y=\frac xn$, $T=L'(1,\chi)M_\chi(x)+O\bigl(\sum_{n\le x}\frac1n\cdot\frac nx(1+\log\frac xn)\bigr)$,
and $\frac1x\sum_{n\le x}(1+\log\frac xn)\le1+\frac1x\bigl(\lfloor x\rfloor\log x-\log\lfloor x\rfloor!\bigr)=O(1)$ by Proposition~\ref{prop:chebyshev}(c). So $L'(1,\chi)M_\chi(x)=\log x+O(1)$.

(b) Since $\log=\Lambda*\mathbf1$, Möbius inversion gives $\Lambda=\mu*\log$, so $\chi(k)\Lambda(k)=\sum_{nm=k}\mu(n)\chi(n)\chi(m)\log m$, and by Lemma~\ref{lem:convsum}
\[
\begin{aligned}S(x,\chi)&=\sum_{n\le x}\frac{\mu(n)\chi(n)}n\sum_{m\le x/n}\frac{\chi(m)\log m}m\eqbr=-L'(1,\chi)M_\chi(x)+O(1),\end{aligned}
\]
with the same error estimate as in (a). By (a), $S(x,\chi)=-\log x+O(1)$. For $\chi_{-4}$ this would mean that $\sum_{p\le x,\,p\equiv1(4)}\frac{\log p}p$ stays bounded: the class $1$ would
receive almost no primes. Note also that (a) forces $L'(1,\chi)\ne0$: a double zero at $1$ is impossible even before Theorem~\ref{thm:lnonzero}.
\end{solution}

\section*{Sources and further reading}

\begin{itemize}
\item P.~G.~L.~Dirichlet, \emph{Beweis des Satzes, dass jede unbegrenzte arithmetische Progression \dots}, Abhandlungen der Königlichen
Preussischen Akademie der Wissenschaften (1837), 45--81.
\item H.~L.~Montgomery and R.~C.~Vaughan, \emph{Multiplicative Number Theory I}, Cambridge, 2007, Chapter~4 (primes in arithmetic
progressions, with Mertens-type estimates).
\item M.~R.~Murty, \emph{Primes in certain arithmetic progressions}, Journal of the Madras University (1988), 161--169 (which progressions admit
a Euclid-style proof).
\item M.~Rubinstein and P.~Sarnak, \emph{Chebyshev's bias}, Experimental Mathematics 3 (1994), 173--197.
\item A.~Granville and G.~Martin, \emph{Prime number races}, American Mathematical Monthly 113 (2006), 1--33. A beautiful survey of the races
between classes.
\item Chebyshev's letter of 1853 to Fuss on the excess of primes $4n+3$ is quoted by Rubinstein and Sarnak (above); the first lead change at $26861$ is J.~Leech, \emph{Note on the distribution of prime numbers}, J.\ London Math.\ Soc.\ 32 (1957), 56--58.
\end{itemize}
